\section*{Chapter \ref{ch:g10:demand} --- Demand}
% figdata: figdata/grade-10/04-demand/out/ (oil_us, oil_changes, oil_summary); Brenland numbers invented.

\begin{solution}{exo:g10:demand:1}
Yes: the quantity falls as the price rises. Between 3 and 5 dalens it falls by $120 - 70 = 50$ loaves for 2 dalens, 25 loaves
per dalen on average (20 between 3 and 4, 30 between 4 and 5).
\end{solution}

\begin{solution}{exo:g10:demand:2}
(a) \omterm{def:g10:demand:substitutes}{Substitutes}. (b) \omterm{def:g10:demand:substitutes}{Complements}. (c) \omterm{def:g10:demand:substitutes}{Substitutes}. (d) \omterm{def:g10:demand:substitutes}{Complements}. (e) \omterm{def:g10:demand:substitutes}{Substitutes}.
\end{solution}

\begin{solution}{exo:g10:demand:3}
At the same prices, higher income goes with less of the cheapest rice: for these households it behaves as an \omterm{def:g10:demand:normal}{inferior good}.
The dearer kinds, bought more when income is higher, behave as \omterm{def:g10:demand:normal}{normal goods}. (A survey across households suggests this; a
change in the income of the same households would show it more firmly.)
\end{solution}

\begin{solution}{exo:g10:demand:4}
The price was highest in 2012 (\usd{4.98} a gallon in 2024 dollars); oil use per person was lowest in 2020 (20.02 barrels).
\end{solution}

\begin{solution}{exo:g10:demand:5}
(a) ``The price of cinema tickets fell, so the \emph{quantity demanded} rose'': a move along the curve. (b) ``Incomes rose, so
\emph{demand} for restaurant meals rose'': at the same prices more is bought, so the curve has shifted.
\end{solution}

\begin{solution}{exo:g10:demand:6}
For $0 \le p \le 5$: $Q = (20 - 4p) + (15 - p) = 35 - 5p$. For $5 < p \le 15$: $Q = 15 - p$. For $p > 15$: $Q = 0$. The kink is at
$p = 5$, where the first buyer stops: $Q = 35 - 25 = 10$ (and $15 - 5 = 10$ on the other piece).
\end{solution}

\begin{solution}{exo:g10:demand:7}
$Q = 3(6 - 2p) = 18 - 6p$ for $p \le 3$. The quantity falls by 6 rides per dalen instead of 2: the slope (in quantity per
dalen) is the sum of the three buyers' slopes. Drawn with the price on the vertical axis, the market line is three times
flatter than one buyer's, and it reaches the price axis at the same point, 3 dalens.
\end{solution}

\begin{solution}{exo:g10:demand:8}
(a) No shift: a move along the curve. (b) Right: driving is dearer, and the bus is a substitute. (c) Left: a new substitute.
(d) Right: more buyers. (e) Right: buyers bring their purchases forward before the rise.
\end{solution}

\begin{solution}{exo:g10:demand:9}
Ana's demand becomes $12 - 2p$ for $p \le 6$. For low fares ($p \le 6$) the \omterm{def:g10:demand:market-demand}{market demand} is $(12 - 2p) + (12 - p) = 24 - 3p$. At a
fare of 3 dalens: $22 - 9 = 13$ rides before, $24 - 9 = 15$ after.
\end{solution}

\begin{solution}{exo:g10:demand:10}
No. The \omterm{def:g10:demand:law}{law of demand} is about a move along one curve, \omterm{def:g10:economic-models:other-things-equal}{other things equal}. In 2025 the shopping centre brought more people
past the cinema and made it easier to reach: the demand curve shifted right. The new point (9 dalens, \num{54000} tickets) lies
on the new curve; on it, a price of 8 dalens would have sold even more tickets. The two observations lie on two different
curves, and comparing them does not test the law.
\end{solution}

\begin{solution}{exo:g10:demand:11}
1980 looks like a move along a demand curve: the price rose and the quantity fell. In 2021 both rose. The economy and travel were
recovering from the first year of the pandemic: incomes and journeys came back, and the demand curve shifted right, pushing up
both the quantity and the price.
\end{solution}

\begin{solution}{exo:g10:demand:12}
A rise in a good's price followed by a rise in the quantity bought, \emph{\omterm{def:g10:economic-models:other-things-equal}{other things equal}}. Before concluding, check that
nothing else moved demand at the same time: buyers' incomes, the prices of \omterm{def:g10:demand:substitutes}{substitutes} and \omterm{def:g10:demand:substitutes}{complements}, expectations of further
price rises (buyers rushing to buy before the next rise), the number of buyers, the quality of the good. The cleanest test is a
change of price that comes from outside the market, or a randomised one, as in the study of the record box.
\end{solution}

\begin{solution}{exo:g10:demand:13}
Between 1978 and 1983 not only the price changed: incomes fell in the recession, cars and heating changed. The line joining the
points mixes the move along the demand curve with the shift of the curve, so it is not the demand curve. If incomes fell at the
same time, demand shifted left, and part of the fall in quantity came from the shift. The move along the true curve, caused by the
price alone, was smaller than the observed fall: the same rise in price with a smaller fall in quantity. With the price on the
vertical axis, the true curve is steeper than the line joining the points.
\end{solution}

\begin{solution}{exo:g10:demand:14}
Rice took most of these households' budget. When its price fell, they became noticeably richer, and they used part of this gain to
replace some rice by meat and other foods they could not afford before, while still getting enough calories. The income effect of
the price fall, towards other foods, outweighed the pull towards the cheaper rice. Two conditions: rice was a very large part of
the budget, and the households were so poor that their diet was mostly rice. Meat is a small part of their budget and a good they
want more of: when its price falls, both effects push them to buy more, so they obey the \omterm{def:g10:demand:law}{law of demand} for meat.
\end{solution}

\begin{solution}{exo:g10:demand:15}
At the old price: poor households buy $1000 \times 2 = 2000$ fewer loaves, richer households $3000 \times 1 = 3000$ more. \omterm{def:g10:demand:market-demand}{Market demand}
at the old price rises by \num{1000} loaves a month: the market curve shifts right, and for the market as a whole bread behaves as a
\omterm{def:g10:demand:normal}{normal good}. The answer would change if there were more poor households, or if they cut their purchases by more: for example, with
\num{2000} poor households the change would be $-4000 + 3000 = -1000$, and bread would behave as \omterm{def:g10:demand:normal}{inferior} for the market.
\end{solution}

\begin{solution}{pb:g10:demand:1}
\textbf{1.} Students at $p = 40$, workers at $p = 60$, retirees at $p = 30$.

\textbf{2.} Students $400 - 200 = 200$, workers $900 - 300 = 600$, retirees $300 - 200 = 100$: 900 passes in all.

\textbf{3.} The workers: 15 passes per dalen, against 10 for the two other groups.

\textbf{4.} At every price the retirees buy 60 more passes after their income rises: the pass is a \omterm{def:g10:demand:normal}{normal good} for them.

\textbf{5.} For $0 \le p \le 30$ all three groups buy: $Q = 400 - 10p + 900 - 15p + 300 - 10p = 1600 - 35p$.

\textbf{6.} For $30 < p \le 40$ the retirees have left: $Q = 1300 - 25p$. For $40 < p \le 60$ only the workers remain: $Q = 900 - 15p$.

\textbf{7.} At $p = 30$ the retirees stop buying: $Q = 1600 - 1050 = 550$ (and $1300 - 750 = 550$). At $p = 40$ the students stop:
$Q = 1300 - 1000 = 300$ (and $900 - 600 = 300$). At each kink the line becomes steeper, since fewer groups respond to the price.

\textbf{8.} Three straight pieces joining $(Q, p) = (1600, 0)$, $(550, 30)$, $(300, 40)$ and $(0, 60)$.

\textbf{9.} 32 lies between 30 and 40: $Q = 1300 - 25 \times 32 = 1300 - 800 = 500$ passes.

\textbf{10.} $900 - 500 = 400$ passes, that is $400/900 \approx 44.4\,\%$ of the passes sold at 20 dalens.

\textbf{11.} Students $400 - 320 = 80$, workers $900 - 480 = 420$, retirees none ($300 - 320 < 0$): the retirees have left the market.
Check: $80 + 420 = 500$.

\textbf{12.} At 20 dalens: $900 \times 20 = \text{\dal{18000}}$ a month. At 32 dalens: $500 \times 32 = \text{\dal{16000}}$. The takings fall by
\dal{2000}.

\textbf{13.} A move along the demand curve: only the price of the pass changed.

\textbf{14.} \omterm{def:g10:demand:substitutes}{Substitutes}: the tram can be used in place of the bus, and its arrival lowers the demand for bus passes.

\textbf{15.} For $0 \le p \le 30$: $Q = (200 - 5p) + (900 - 15p) + (300 - 10p) = 1400 - 30p$. For $30 < p \le 40$: $Q = (200 - 5p) +
(900 - 15p) = 1100 - 20p$.

\textbf{16.} At 20 dalens: $1400 - 600 = 800$. At 32 dalens: $1100 - 640 = 460$.

\textbf{17.} $800 - 460 = 340$ passes, that is $340/800 = 42.5\,\%$.

\textbf{18.} The students now respond to the price by 5 passes per dalen instead of 10: from 20 to 32 dalens they give up $100 - 40 = 60$
passes instead of $200 - 80 = 120$. The workers and retirees lose the same numbers as before, so the total loss is 60 smaller. Part of the
students who would have left at a higher fare have already left for the tram.

\textbf{19.} $900 - 800 = 100$ passes: a shift of the demand curve, caused by the tram, at an unchanged price.

\textbf{20.} The fare rise from 20 to 32 dalens cuts the passes sold by 400 a month without the tram line ($44.4\,\%$) and by 340 with it
($42.5\,\%$).
\end{solution}
